\chapter{Background and Related Work}
\label{ch:background}

A.L.F.R.E.D.\ sits at the intersection of four lines of work: small on-device
language models, language-model tool use and agents, model routing and cascading,
and the distillation of agent experience into reusable skills. This chapter
reviews each in turn and ends by stating precisely what gap the thesis fills. A
conceptual thread runs underneath the whole design and is worth naming first. The
two-tier split, a fast automatic path and a slow deliberate one, is the
computational analogue of two old ideas: the reflex arc, a stimulus-driven motor
response routed without higher deliberation~\citep{sherrington1906}, and the
dual-process account of cognition, in which a fast automatic system handles the
routine and a slow effortful system handles the novel~\citep{kahneman2011}. The
reflexer is the fast path; the thinker is the slow one.

\section{Small on-device language models}
\label{sec:bg-slm}

The premise that a small model is sufficient for most agentic work is now argued
directly. \citet{belcak2025slm} contend that small language models, not frontier
models, are the right substrate for agentic systems, because agent workloads are
dominated by narrow, repetitive sub-tasks that a small specialized model can serve
at a fraction of the cost. A.L.F.R.E.D.\ is a concrete instantiation of that
thesis for the on-device action setting, and supplies the mechanism (distilled
patterns) that the position paper leaves open.

On the systems side, TinyAgent~\citep{erdogan2024tinyagent} demonstrates reliable
function calling at the edge with small models, using fine-tuning and a
constrained decoder to keep a small model on the rails of a fixed tool schema.
A.L.F.R.E.D.\ shares the goal of reliable small-model tool use on-device but takes
a different route to reliability. Rather than fine-tuning the model to a schema,
it hands the model a distilled pattern at inference time, keeping the model itself
untouched and the capability set extensible by adding patterns. Industry on-device
stacks (Gemini Nano through Android AICore, and Apple Intelligence with Private
Cloud Compute for overflow) make the same architectural bet of a local small model
with selective escalation; A.L.F.R.E.D.\ is a research-grade realization of that
pattern whose tiers can be measured in isolation.

\section{Tool use and agent reasoning}
\label{sec:bg-agents}

The thinker tier implements the now-standard reasoning-and-acting loop introduced
by ReAct~\citep{yao2022react}, in which a model interleaves reasoning steps with
tool calls and observations until a task is complete. The broader line on teaching
models to use tools includes Toolformer~\citep{schick2023toolformer}, which shows a
model can learn to call APIs from self-supervised data, and
Gorilla~\citep{patil2023gorilla}, which connects a model to massive real API
collections and confronts the problem of selecting the correct call from many
candidates. These establish that capable models can use tools; the question this
thesis asks is the inverse one, namely how little model is needed once the
\emph{which-tool} and \emph{which-shape} decisions have been settled in advance by
a distilled pattern. Where ReAct spends model capability at run time to decide how
to act, the reflexer spends it once, offline, and replays the decision.

\section{Routing and cascading}
\label{sec:bg-routing}

Deciding which model or path should serve a request is the routing problem.
RouteLLM~\citep{ong2024routellm} learns to route between a strong and a weak model
from human preference data, optimizing the cost-quality trade-off;
\citet{dekoninck2024unifiedrouting} give a unified treatment of routing and
cascading and frame multi-tier systems as structured cascades, which is exactly
how A.L.F.R.E.D.'s router-then-reflexer-or-thinker path should be understood. The
difference in emphasis matters for this thesis. The routing literature largely
treats the downstream models as black boxes and optimizes the \emph{decision};
A.L.F.R.E.D.\ instead invests in making the cheap downstream path itself reliable
through distillation, and then reports honestly that the \emph{router} remains the
binding constraint on end-to-end accuracy (Section~\ref{sec:routing}). The two
concerns are complementary: a better router and a more reliable reflexer multiply
rather than substitute.

\section{Skill memory and procedural distillation}
\label{sec:bg-skills}

The closest neighbours come from work that turns agent experience into reusable
skills. Voyager~\citep{wang2023voyager} introduced the skill-library framing, in
which an agent accumulates a growing library of executable skills and composes
them; Memp~\citep{fang2025memp} studies procedural memory directly, abstracting
successful trajectories into reusable scripts. A.L.F.R.E.D.'s patterns can be read
as the deterministic compilation of exactly such procedural abstractions: where
Memp keeps scripts the agent re-interprets, A.L.F.R.E.D.\ freezes them into a
template whose only run-time freedom is slot-filling.

The Memento line is the nearest published relative. Memento~\citep{zhou2025memento}
adapts an agent without updating model weights by retrieving from a growing case
bank, a non-parametric personalization that is the same high-level strategy
A.L.F.R.E.D.\ uses with patterns. Its successor,
Memento-Skills~\citep{zhou2026mementoskills}, adds a \emph{learned} skill-router
and reflective read-write learning over free-form markdown skills. A.L.F.R.E.D.\
contrasts on two axes: its skills are \emph{compiled, deterministic} patterns
rather than free-form documents, and (as the evaluation will argue) the production
end is best served by a constrained executor whose command shape is fixed, trading
the generality of a reflective skill for the reliability of a frozen one. The
opposite design point, updating model weights to personalize an agent, is
represented by reinforcement-learning approaches such as
OpenClaw-RL~\citep{wang2026openclawrl}; this is the path A.L.F.R.E.D.\ explicitly
does not take, because weight updates sacrifice the plug-and-play, inspectable
extensibility that a symbolic pattern provides.

A note on naming is needed to avoid confusion. Our \emph{reflexer} is unrelated to
Reflexion~\citep{shinn2023reflexion}, which uses verbal self-feedback as a form of
reinforcement to improve an agent across attempts. Reflexion adds deliberation; the
reflexer removes it. The shared root word denotes opposite design intents.

\section{Knowledge distillation, classical and symbolic}
\label{sec:bg-distillation}

The mechanism that makes the cheap tier work is distillation, and its lineage is
explicit. Classical knowledge distillation compresses a large teacher network into
the \emph{weights} of a smaller student, transferring the teacher's soft behaviour
into a deployable model~\citep{hinton2015distilling}. A.L.F.R.E.D.\ distills along
the same teacher-to-student axis but changes the target: the teacher's knowledge of
device conventions is compressed into an inspectable \emph{symbolic} artifact, a
pattern, rather than into student parameters. The trade is deliberate. Weight
distillation yields a general but opaque student that must be retrained to extend;
symbolic distillation yields a student that is plug-and-play, auditable, and
composable, at the cost of covering only distilled intents. This positions the
contribution against both the distillation literature (same teacher-student idea,
different artifact) and the programmatic-prompt line such as
DSPy~\citep{khattab2024dspy}, which compiles declarative model calls into optimized
pipelines; A.L.F.R.E.D.\ similarly compiles intents into fixed command shapes, but
targets on-device execution rather than prompt-program optimization.

\section{Evaluation of agents}
\label{sec:bg-eval}

Agent evaluation has converged on task suites that score completion of realistic
multi-step goals: AgentBench~\citep{liu2023agentbench} across diverse environments,
and AssistantBench~\citep{yoran2024assistantbench} for realistic, time-consuming
web tasks. A.L.F.R.E.D.'s benchmarks share their realism goal but differ in
instrument. Rather than scoring trajectory completion, the settings benchmarks
score the \emph{end state} of a mutable backend with a no-collateral constraint
(Chapter~\ref{ch:method}), which isolates the specific failure of emitting a
plausible command that produces the wrong device state. The skill interface itself
follows the emerging \texttt{SKILL.md} convention for equipping agents with
packaged capabilities~\citep{anthropic2025skills}, with which A.L.F.R.E.D.\ is
structurally compatible.

\section{The gap}
\label{sec:bg-gap}

Across these lines, three gaps remain that this thesis addresses. First, the
small-model and routing literatures optimize the \emph{decision} of which model to
use but rarely quantify how reliable the cheap path can be made, or by what
mechanism; A.L.F.R.E.D.\ measures exactly this and finds distillation, not scale,
to be the lever (Chapter~\ref{ch:eval}). Second, the skill-memory line keeps
skills as re-interpreted scripts or free-form documents, and the consequence of
\emph{freezing} them into deterministic templates, higher reliability at the cost
of generality, has not been measured against a strong free-form baseline on a
realistic device API. Third, and most specific, no prior work separates the two
difficulties of device action, \emph{discovery} (which key) and \emph{convention}
(what the value means), then shows that model scale fixes the first but not the
second while a distilled pattern fixes both. That separation, and the
teacher-quality result that follows from it, is the contribution the remaining
chapters defend.
